
\subsubsection{Dataset}

We use Waterbirds, a bird classification dataset with 11,788 images across four groups: landbirds on land (3,498 samples), landbirds on water (184 samples), waterbirds on land (56 samples), and waterbirds on water (1,057 samples). The dataset contains a spurious correlation between background type and bird type, enabling models to achieve high average accuracy while failing on minority groups where the correlation does not hold.

\subsubsection{Model and Training}

We use ResNet-50 pretrained on ImageNet and fine-tune on Waterbirds. Training uses SGD with learning rate 0.001, momentum 0.9, and batch size 128. Standard data augmentation includes random cropping and horizontal flipping. For ERM, we train with cross-entropy loss for 67 epochs. For Group-DRO, we train with worst-group loss minimization for 45 epochs.

We conduct the primary experiment (geometric signature comparison) with one random seed. The stochastic gradient descent trajectory analysis uses three seeds (42, 43, 44). Interpolation path experiments train for only 5 epochs, substantially less than the full training protocol.

\subsubsection{Baselines}

Standard ERM minimizes average cross-entropy loss. Group-DRO (Sagawa et al., 2020) minimizes worst-group loss via adaptive group reweighting, requiring group labels during training. We compare Hessian eigenvalue structure and alignment metrics between converged ERM and Group-DRO solutions.

\subsubsection{Evaluation Metrics}

The primary metric is curvature subspace alignment A(θ). Secondary metrics include worst-group accuracy (minimum accuracy across the four subgroups), outlier eigenvalue count (number exceeding Marchenko-Pastur bulk edge), and directional bias during training (difference between bulk and outlier alignment).
